\documentclass[12pt]{article}
\usepackage[T1]{fontenc}
\usepackage[utf8]{inputenc}
\usepackage{lmodern}
\usepackage{textcomp}
\usepackage{enumitem}
\usepackage{geometry}
\geometry{margin=1in}
\begin{document}
\begin{center}
Answer Key - Psychology - Graduate Level Assessment 16 \
\small (Answers are ordered exactly as in the question paper.)
\end{center}

\section*{Multiple Choice}
\begin{enumerate}[label=\arabic*.]
\item \textbf{Answer:} A
\\ \textit{Options:} A) "Genetic blueprint" hypothesis; B) "Innate potential" hypothesis; C) Environmental impact hypothesis; D) "Experience-driven" development theory; E) "Blank slate" hypothesis
\\ \textit{Note:} Correct option: A - "Genetic blueprint" hypothesis
\item \textbf{Answer:} C
\\ \textit{Options:} A) too little curvature of the cornea and lens; B) too much curvature of the retina and lens; C) too much curvature of the cornea and lens; D) too much curvature of the iris and cornea; E) incorrect alignment of the cornea and retina
\\ \textit{Note:} Correct option: C - too much curvature of the cornea and lens
\item \textbf{Answer:} A
\\ \textit{Options:} A) By observing individual's behaviors and actions; B) Through the analysis of social media usage and engagement patterns; C) By analyzing the country's economic factors; D) By examining the literacy and numeracy rates; E) Through standardized achievement tests in educational institutions
\\ \textit{Note:} Correct option: A - By observing individual's behaviors and actions
\item \textbf{Answer:} A
\\ \textit{Options:} A) Recovery rate in individuals using self-help resources; B) Incidence in the population of the disorder being eared; C) Recovery rate in patients using alternative therapies; D) Improvement rate in patients receiving no therapy; E) Cure rate for older, outdated treatment methods
\\ \textit{Note:} Correct option: A - Recovery rate in individuals using self-help resources
\item \textbf{Answer:} B
\\ \textit{Options:} A) Recognize faces; B) move his left hand; C) Hear sounds correctly; D) understand information he reads; E) speak intelligibly
\\ \textit{Note:} Correct option: B - move his left hand
\end{enumerate}

\section*{True / False}
\begin{enumerate}[label=\arabic*.]
\setcounter{enumi}{5}
\item \textbf{Answer:} False
\\ \textit{Options:} A) True; B) False
\\ \textit{Note:} LLM-generated false statement. Correct answer: 'Simultaneous processing'.
\item \textbf{Answer:} False
\\ \textit{Options:} A) True; B) False
\\ \textit{Note:} LLM-generated false statement. Correct answer: 'a conditioned stimulus'.
\item \textbf{Answer:} False
\\ \textit{Options:} A) True; B) False
\\ \textit{Note:} LLM-generated false statement. Correct answer: 'opportunities for catharsis'.
\item \textbf{Answer:} True
\\ \textit{Options:} A) True; B) False
\\ \textit{Note:} LLM-generated true statement based on MCQ.
\item \textbf{Answer:} False
\\ \textit{Options:} A) True; B) False
\\ \textit{Note:} LLM-generated false statement. Correct answer: 'delay between the CS and the UCS is too long'.
\end{enumerate}

\section*{Long Form Response}
\begin{enumerate}[label=\arabic*.]
\setcounter{enumi}{10}
\item \textbf{Answer:} pleasure and a visual stimulus. Explain why this is the correct answer and provide supporting evidence. Reference key facts or concepts that support this choice.
\\ \textit{Note:} Long-form derived from MCQ; award credit for stating the correct idea and giving supporting reasoning.
\item \textbf{Answer:} Neurotic symptoms are a result of childhood trauma.. Explain why this is the correct answer and provide supporting evidence. Reference key facts or concepts that support this choice.
\\ \textit{Note:} Long-form derived from MCQ; award credit for stating the correct idea and giving supporting reasoning.
\end{enumerate}

\vfill
\noindent\textit{End of Answer Key}
\end{document}